\section{Related work}
\label{sec:related}

The use of an attention mechanism in \ac{nlp} and in particular \ac{nmt} can be traced back to \citet{bahdanau15alignandtranslate} and \citet{cho2014learning}, and most contemporaneous implementations are based on the formulation from \citet{luong2015effective}. Attention was shortly adapted (successfully) to other \ac{nlp} tasks, often achieving then state-of-the-art performance in reading comprehension \citep{cheng-etal-2016-long}, natural language inference \citep{parikh-etal-2016-decomposable} or abstractive summarization \citep{paulus2017deep} to cite a few. Multi-headed attention was first introduced by \citet{vaswani2017attention} for \ac{nmt} and English constituency parsing, and later adopted for transfer learning \citep{radford2018improving,devlin2018bert}, language modeling \citep{dai2019transformer,radford2019language}, or semantic role labeling \citep{strubell2018linguistically}, among others.

There is a rich literature on pruning trained neural networks, going back to \citet{lecun1989optimal} and \citet{hassibi92second} in the early 90s and reinvigorated after the advent of deep learning, with two orthogonal approaches: fine-grained ``weight-by-weight'' pruning \citep{han2015learning} and structured pruning \citep{Anwar:2017:SPD:3051701.3005348,li2016pruning,molchanov2016pruning}, wherein entire parts of the model are pruned. In \ac{nlp} , structured pruning for auto-sizing feed-forward language models was first investigated by \citet{murray-chiang-2015-auto}. More recently, fine-grained pruning approaches have been popularized by \citet{see-etal-2016-compression} and \citet{kim-rush-2016-sequence} (mostly on \ac{nmt}).

Concurrently to our own work, \citet{voita2019analyzing} have made to a similar observation on multi-head attention. Their approach involves using LRP \citep{binder2016layer} for determining important heads and looking at specific properties such as attending to adjacent positions, rare words or syntactically related words. They propose an alternate pruning mechanism based on doing gradient descent on the mask variables $\xi_h$. While their approach and results are complementary to this paper, our study provides additional evidence of this phenomenon beyond \ac{nmt}, as well as an analysis of the training dynamics of pruning attention heads.